\section{Introduction}

High-order finite and spectral element methods are widely used in scientific simulations across many domains, including computational fluid dynamics, electromagnetics, seismology, and magnetohydrodynamics. These methods divide the physical domain into an unstructured mesh of elements and sample each physical variable at multiple points within every element. The resulting numerical fields combine irregular connectivity among elements with a common arrangement of points within each element. We refer to this data type as \emph{high-order finite and spectral element data}, or \emph{high-order element data} for short.

High-order element simulations can generate large data volumes. In one NekRS~\cite{fischer2022nekrs} simulation conducted by our collaborators, the output at a single time step is approximately 6~TB. Storing and transferring data at this scale can therefore incur substantial costs.

Error-bounded lossy compression~\cite{liang2018error,liu2024high,zhao2021optimizing,liang2023sz3,lindstrom2014fixed,li2023lossy,di2025survey} can substantially reduce scientific data volumes while ensuring that the error of every reconstructed value remains within a user-specified bound. This combination of high compression ratios and explicit error control makes it well suited to storing and transferring high-order element data. In practical data-management workflows, compression and decompression must also be sufficiently fast so that their computational overhead does not offset the savings in storage and data movement. These requirements motivate us to investigate error-bounded lossy compression specialized for high-order element data.

General-purpose compressors can process high-order element data, but they do not fully account for the organization of points within and across elements. Compressors designed for multidimensional arrays, including SZ2~\cite{tao2017significantly}, SZ3~\cite{liang2023sz3,zhao2021optimizing}, QoZ~\cite{liu2024high}, ZFP~\cite{lindstrom2014fixed}, SPERR~\cite{li2023lossy}, and MGARD~\cite{gong2023mgard}, can be used through their official APIs with only simple preprocessing. However, they are not designed for the different statistical patterns associated with element-local positions, and all except MGARD assume uniform spacing. Compressors designed for unstructured data build predictions based on spatial proximity or mesh connectivity~\cite{wu2025enabling,gong2024general,ren2024prediction}. Although these predictors capture relationships in physical space, they do not exploit the quasi-structured arrangement of points within each element or the correspondence among local positions across elements.

Compression approaches have also been studied specifically for high-order element data. Shah et al.~\cite{shah2026solver} integrate SZ3 into the NekRS parallel output path to improve I/O performance, but do not modify SZ3's compression model to improve its effectiveness for this data organization. Another approach exploits the particular structure of the data~\cite{marin2016spectral}, but it does not provide a strict pointwise error guarantee, making it difficult to ensure the reliability of reconstructed data in downstream analysis.

We address these limitations with two complementary error-bounded compressors. The first, FDWZ, targets high throughput while retaining effective storage reduction. Building on the multilevel interpolation framework of SZ3~\cite{liang2023sz3,zhao2021optimizing}, FDWZ accounts for the nonuniform spacing of points within each element and uses multiple Huffman trees for data groups with different distributions. The second, MRCZ, targets maximum storage reduction. It represents each element using transform coefficients and residuals, then groups corresponding coefficients across elements to improve compression. This representation requires more computation than FDWZ but produces substantially smaller compressed data.

We evaluate FDWZ and MRCZ on high-order element data from NekRS, Nektar++, and MFEM and compare them with SZ3, SZ2, QoZ, ZFP, SPERR, and MGARD using each baseline's official API. The evaluation covers three concerns for scientific data management: storage reduction at comparable data quality, compression and decompression throughput, and end-to-end movement of compressed fields between facilities. At matched PSNRs, FDWZ achieves a geometric-mean compression-ratio improvement of 73.1\% over the best-performing baseline in our large-scale experiments. Across 40 throughput cases spanning lower and higher CPU allocations, FDWZ ranks among the two fastest methods in 38 cases, including first in 18 of 20 cases at the lower core counts. In our data-transfer experiments using Globus~\cite{chard2017globus}, FDWZ matches or outperforms the fastest baseline, reducing end-to-end time by up to 19.2\%. At matched PSNRs in the large-scale experiments, MRCZ improves the compression ratio over FDWZ by a geometric mean of 108.5\%. These results show that FDWZ provides a practical balance between speed and storage reduction, while MRCZ is suitable when compression ratio is the main concern.

The main contributions of this work are as follows:
\begin{itemize}
\item We develop FDWZ, which combines distance-weighted multilevel interpolation and multi-tree Huffman encoding to improve compression ratio while maintaining high throughput under a strict pointwise error bound.
\item We develop MRCZ, which combines an element-wise coefficient-and-residual representation with a mode-major coefficient layout across elements to achieve substantially higher compression ratios than FDWZ at the same data distortion while preserving a strict pointwise error bound.
\item We build a common partitioned processing workflow for FDWZ, MRCZ, and six representative baselines and use it to evaluate storage reduction, reconstruction quality, codec throughput, and end-to-end cross-facility data-movement performance across eight datasets.
\end{itemize}

The rest of this paper is organized as follows. Section~\ref{sec:related work} reviews compression methods for structured, unstructured, and high-order element data. Section~\ref{sec:problem formulation} describes the data characteristics and formulates the compression problem. Section~\ref{sec:design overview} presents the parallel execution model and the element-local design of FDWZ and MRCZ. Sections~\ref{sec:SZ3-HOHS} and~\ref{sec:mrcz} describe FDWZ and MRCZ, respectively. Section~\ref{sec:performance evaluation} presents the experimental evaluation, and Section~\ref{sec:conclusion} concludes the paper.
